\documentclass[11pt,a4paper]{article}

\usepackage[margin=2.5cm]{geometry}
\usepackage{amsmath,amssymb,amsthm}
\usepackage{bm}
\usepackage{booktabs}
\usepackage{xcolor}
\usepackage{hyperref}
\usepackage{enumitem}

\newtheorem{lemma}{Lemma}[section]
\newtheorem{theorem}[lemma]{Theorem}
\newtheorem{remark}[lemma]{Remark}

\newcommand{\R}{\mathbb{R}}
\newcommand{\dom}{\Omega}
\newcommand{\pore}[1]{\Omega_{#1}}
\newcommand{\iface}{\gamma}
\newcommand{\vel}{\bm{u}}
\newcommand{\pres}{p}
\newcommand{\stress}{\bm{\sigma}}
\newcommand{\normal}{\bm{n}}
\newcommand{\trac}{\bm{\lambda}}
\newcommand{\DtN}{\mathbf{G}}
\newcommand{\Schur}{\mathbf{S}}
\newcommand{\one}{\mathbf{1}}
\newcommand{\zero}{\bm{0}}

\newcommand{\modeset}{\mathcal{M}}
\newcommand{\loadvec}{\bm{f}}
\newcommand{\respvel}{\mathbf{U}}
\newcommand{\resppres}{\mathbf{X}}

\hypersetup{colorlinks=true,linkcolor=blue,urlcolor=blue,citecolor=blue}

\begin{document}

\title{\textbf{Mathematical Properties of the Affine-DDPNM}}
\author{}
\date{\today}
\maketitle

\section{Preliminaries and notation}

We adopt the domain partition, finite-element discretisation, and notation
of the accompanying methodology document.  The pore domain
$\dom \subset \R^3$ is partitioned into $N$ non-overlapping subdomains
$\{\pore{i}\}_{i=1}^N$ with $m$ internal interfaces
$\{\iface_k\}_{k=1}^m$, assumed planar.
Subdomain $\pore{i}$ has $m_i$ ports, of which $m_i^u$ are internal
(unknown-traction) and $m_i^k$ are boundary (known-traction).
Each $\pore{i}$ is required to contain at least one positive-measure
no-slip wall segment, ensuring that the local Stokes operator with
mixed boundary conditions is invertible (no rigid-body modes).

\medskip\noindent
\textbf{Interface traction basis.}
On each interface $\iface$, we introduce a local orthonormal frame
$\{\normal_\iface, \bm{t}_{\iface,1}, \bm{t}_{\iface,2}\}$ with origin at the
interface centroid $\bm{c}_\iface$, and dimensionless in-plane coordinates
\begin{equation}\label{eq:local-coords}
  s(\bm{x}) = \frac{(\bm{x}-\bm{c}_\iface)\cdot\bm{t}_{\iface,1}}{a_\iface},\qquad
  t(\bm{x}) = \frac{(\bm{x}-\bm{c}_\iface)\cdot\bm{t}_{\iface,2}}{b_\iface},
\end{equation}
where $a_\iface, b_\iface$ are the half-spans of $\iface$ along the two
tangent directions, ensuring $s,t = O(1)$ on $\iface$ under mild
shape-regularity.  The \emph{mode index set} is
\begin{equation}\label{eq:modeset}
  \modeset = \{0,s,t\} \times \{\normal, \bm{t}_1, \bm{t}_2\},
\end{equation}
with scalar shape functions $\varphi_0 = 1$, $\varphi_s = s$, $\varphi_t = t$.
For a subdomain $\pore{i}$ adjacent to $\iface$, define the \emph{signed}
direction vectors
\begin{equation}\label{eq:signed-dir}
  \bm{d}_{i,\iface,\normal} = \normal_{i,\iface},\qquad
  \bm{d}_{i,\iface,\bm{t}_\nu} = \sigma_{i,\iface}\,\bm{t}_{\iface,\nu}
  \;\;(\nu=1,2),
\end{equation}
where $\sigma_{i,\iface}=+1$ if $\pore{i}$ is the first member of the
interface pair and $\sigma_{i,\iface}=-1$ otherwise, satisfying
$\bm{d}_{j,\iface,\beta} = -\bm{d}_{i,\iface,\beta}$ for the two pores
$i,j$ sharing $\iface$.  Throughout we abbreviate $\bm{d}_{i,\gamma,\beta}$.

The \emph{pseudostress} $\stress(\vel,\pres) := -\pres\mathbf{I}
+ \nu\nabla\vel$ is used throughout.  On each interface $\iface$, the
applied traction is expanded as
\begin{equation}\label{eq:traction}
  -\stress(\vel,\pres)\,\normal_{i,\iface}
  = \sum_{(\alpha,\beta)\in\modeset}
    \lambda_{\iface,\alpha,\beta}\;
    \varphi_\alpha(s,t)\,\bm{d}_{i,\iface,\beta},
\end{equation}
where the minus sign on the left follows the classic convention.
The nine coefficients $\bm{\lambda}_\iface
= (\lambda_{\iface,0,\normal},\dots,\lambda_{\iface,t,\bm{t}_2})
\in \R^9$ are the unknown interface degrees of freedom.

For a subdomain $\pore{i}$, let $n_i = 9m_i^u + m_i^k$ be the total number of
local traction unknowns and $\trac_i \in \R^{n_i}$ the vector of all local
coefficients.  The Boolean restriction maps $R_i^u \in \{0,1\}^{n_i^u \times
9m}$ and $R_i^k \in \{0,1\}^{m_i^k \times m^k}$ and the global unknown
$\trac \in \R^{9m}$ are defined with block size $9\times9$ per interface.

\medskip\noindent
\textbf{Unit-response loads.}
For each mode $(\alpha,\beta) \in \modeset$ on an internal interface
$\iface_{i,r}$ of $\pore{i}$, the Neumann load vector
$\loadvec_{i,r}^{(\alpha,\beta)} \in \R^{n_{i,u}}$ has entries
\begin{equation}\label{eq:load}
  \bigl[\loadvec_{i,r}^{(\alpha,\beta)}\bigr]_j
  = \int_{\iface_{i,r}} \varphi_\alpha(s,t)\;
    \bm{d}_{i,\gamma,\beta} \cdot \bm{\phi}_j^{(i)} \;\mathrm{d}S,
\end{equation}
where $\{\bm{\phi}_j^{(i)}\}$ is the local velocity basis.
By the sign convention~\eqref{eq:traction}, a positive coefficient
$\lambda$ produces a traction directed \emph{into} the pore; the load
vector carries no extra minus sign---it is the right-hand side of the
discrete momentum equation.  For boundary ports, only the mode
$(\alpha,\beta) = (0,\normal)$ is used with a \emph{prescribed} coefficient
equal to the inlet/outlet pressure.

\medskip\noindent
\textbf{Local Stokes system.}
For each unit load, the discrete Taylor--Hood $[P_2]^3$--$P_1$ system reads
\begin{equation}\label{eq:local-sys}
  \begin{bmatrix} A_i & B_i^\top \\ B_i & -\varepsilon M_{p,i} \end{bmatrix}
  \begin{bmatrix} \respvel_{i,r}^{(\alpha,\beta)} \\
                   \resppres_{i,r}^{(\alpha,\beta)} \end{bmatrix}
  = \begin{bmatrix} \loadvec_{i,r}^{(\alpha,\beta)} \\ \bm{0} \end{bmatrix},
\end{equation}
where $A_i$ is SPD, $B_i$ is the divergence matrix,
$\varepsilon \ge 0$ is an optional pressure stabilisation parameter, and
$M_{p,i}$ is the pressure mass matrix.  In what follows we set
$\varepsilon = 0$ (no stabilisation), so that exact discrete
incompressibility $B_i\respvel_i = \bm{0}$ holds.  All numerical results
reported use $\varepsilon = 0$; for $\varepsilon > 0$ the mass conservation
statement in Theorem~\ref{thm:main} requires a modification.

\medskip\noindent
\textbf{Local compliance map.}
The local Neumann-to-Dirichlet (compliance) map
$\DtN_i \in \R^{n_i \times n_i}$ has entries
\begin{equation}\label{eq:dtn}
  \DtN_i\bigl[(\gamma,\alpha,\beta),\,(\gamma',\alpha',\beta')\bigr]
  = \bigl[\loadvec_{i}^{(\alpha,\beta)}(\gamma)\bigr]^\top\,
    \respvel_{i}^{(\alpha',\beta')}(\gamma'),
\end{equation}
where the index triple $(\gamma,\alpha,\beta)$ ranges over the $n_i$ modes.
$\DtN_i$ is symmetric positive semi-definite (proved in Lemma~\ref{lem:sym}).
It is partitioned as
\[
  \DtN_i =
  \begin{bmatrix}
    \DtN_i^{uu} & \DtN_i^{uk} \\[2pt]
    \DtN_i^{ku} & \DtN_i^{kk}
  \end{bmatrix},
\]
where $\DtN_i^{uu} \in \R^{9m_i^u \times 9m_i^u}$ maps unknown tractions,
$\DtN_i^{uk} \in \R^{9m_i^u \times m_i^k}$ couples unknown to known, and
$\DtN_i^{kk} \in \R^{m_i^k \times m_i^k}$ involves boundary ports only.

\medskip\noindent
\textbf{Weighted velocity moments.}
For each mode $(\alpha,\beta)$ on interface $\iface_{i,r}$, define
\begin{equation}\label{eq:moment}
  M_{i,r}^{(\alpha,\beta)}
  = \int_{\iface_{i,r}} \varphi_\alpha(s,t)\;
    \vel_{i,h} \cdot \bm{d}_{i,\gamma,\beta} \;\mathrm{d}S.
\end{equation}
Only $M_{i,r}^{(0,\normal)} = \int_\iface \vel\cdot\normal\,\mathrm{d}S$
is the volumetric flow rate.  The remaining eight quantities are weighted
velocity moments (linear normal flux moments $M^{(s,\normal)},M^{(t,\normal)}$
and tangential velocity moments $M^{(\alpha,\bm{t}_\nu)}$).
From $\bm{d}_{j,\gamma,\beta} = -\bm{d}_{i,\gamma,\beta}$ we obtain the
consistency condition
$M_{i,r}^{(\alpha,\beta)} + M_{j,r}^{(\alpha,\beta)} = 0$ for any
continuous velocity field across $\iface_r$.

\medskip\noindent
\textbf{Global Schur complement.}
Enforcing continuity of all nine weighted velocity moments on every internal
interface yields the global system
\begin{equation}\label{eq:schur}
  \Schur\,\trac = \bm{F},\qquad
  \Schur = \sum_{i=1}^N (R_i^u)^\top \DtN_i^{uu} R_i^u,\qquad
  \bm{F} = -\sum_{i=1}^N (R_i^u)^\top \DtN_i^{uk} R_i^k \trac^k,
\end{equation}
where $\trac^k \in \R^{m^k}$ collects prescribed boundary tractions.
$\Schur \in \R^{9m \times 9m}$ is symmetric; its positive definiteness
is established in Theorem~\ref{thm:main} under an anchoring condition.

% ======================================================================
\section{Mathematical properties}
% ======================================================================

\begin{lemma}[Local well-posedness]\label{lem:wellposed}
  Assume each $\pore{i}$ contains a positive-measure no-slip wall segment.
  Then for every interface-traction vector $\trac_i \in \R^{n_i}$, the
  discrete Stokes subproblem~\eqref{eq:local-sys} in $\pore{i}$ with
  $\varepsilon = 0$ admits a unique solution
  $(\respvel_i, \resppres_i)$.
\end{lemma}

\begin{proof}
  The matrix $A_i$ is SPD.  With the wall Dirichlet condition, the
  Taylor--Hood pair satisfies the discrete inf--sup condition on
  $\pore{i}$, so $B_i$ has full row rank and the Schur complement
  $B_i A_i^{-1} B_i^\top$ is SPD.  Hence the saddle-point matrix
  $K_i$ is invertible.  The right-hand side of~\eqref{eq:local-sys}
  is a linear combination of the primitive load vectors; a unique
  solution exists for any $\trac_i$.
\end{proof}

% ----------------------------------------------------------------------

\begin{lemma}[Discrete incompressibility of unit responses]\label{lem:incomp}
  For each mode $(\alpha,\beta) \in \modeset$ on an internal interface
  $\iface_{i,r}$, the response velocity satisfies
  $B_i\,\respvel_{i,r}^{(\alpha,\beta)} = \bm{0}$.
\end{lemma}

\begin{proof}
  Immediate from the second block row of~\eqref{eq:local-sys} with
  $\varepsilon = 0$.  For reference we also give the explicit form:
  $\respvel_{i,r}^{(\alpha,\beta)} = H_i \loadvec_{i,r}^{(\alpha,\beta)}$
  where
  \begin{equation}\label{eq:Hdef}
    H_i := A_i^{-1} - A_i^{-1} B_i^\top M_i^{-1} B_i A_i^{-1},\qquad
    M_i := B_i A_i^{-1} B_i^\top,
  \end{equation}
  and $B_i H_i = 0$ is verified directly.
\end{proof}

% ----------------------------------------------------------------------

\begin{lemma}[Symmetry and energy identity]\label{lem:sym}
  $\DtN_i$ defined by~\eqref{eq:dtn} is symmetric positive semi-definite.
  For any $\trac_i \in \R^{n_i}$,
  \begin{equation}\label{eq:energy}
    \trac_i^\top \DtN_i \trac_i
    = \respvel_i^\top A_i \respvel_i \ge 0,
  \end{equation}
  where $\respvel_i = \sum_{\gamma,\alpha,\beta}
    \lambda_{i,\gamma,\alpha,\beta}\,
    \respvel_{i}^{(\alpha,\beta)}(\gamma)$
  is the velocity field induced by $\trac_i$.
\end{lemma}

\begin{proof}
  Let $\mu = (\gamma,\alpha,\beta)$, $\mu' = (\gamma',\alpha',\beta')$.
  From Lemma~\ref{lem:incomp} and~\eqref{eq:Hdef},
  $\respvel_i^{(\mu)} = H_i \loadvec_i^{(\mu)}$.
  The identity $H_i^\top A_i H_i = H_i$ follows from direct expansion:
  \begin{align*}
    H_i^\top A_i H_i
    &= (A_i^{-1} - A_i^{-1}B_i^\top M_i^{-1}B_i A_i^{-1})\,A_i\,
       (A_i^{-1} - A_i^{-1}B_i^\top M_i^{-1}B_i A_i^{-1}) \\
    &= A_i^{-1} - 2A_i^{-1}B_i^\top M_i^{-1}B_i A_i^{-1}
       + A_i^{-1}B_i^\top M_i^{-1}(B_i A_i^{-1}B_i^\top)M_i^{-1}B_i A_i^{-1}\\
    &= A_i^{-1} - A_i^{-1}B_i^\top M_i^{-1}B_i A_i^{-1}
     = H_i,
  \end{align*}
  since $B_i A_i^{-1} B_i^\top = M_i$.  Now:
  \begin{align*}
    \DtN_i(\mu,\mu')
    &= \bigl[\loadvec_i^{(\mu)}\bigr]^\top \respvel_i^{(\mu')}
     = \bigl[\loadvec_i^{(\mu)}\bigr]^\top H_i \loadvec_i^{(\mu')} \\
    &= \bigl[\loadvec_i^{(\mu)}\bigr]^\top (H_i^\top A_i H_i)
       \loadvec_i^{(\mu')}
     = \bigl[\respvel_i^{(\mu)}\bigr]^\top A_i \respvel_i^{(\mu')}.
  \end{align*}
  Symmetry follows from symmetry of $A_i$.
  Summing over $\mu,\mu'$ with coefficients $\lambda_{i,\mu}$ yields
  $\trac_i^\top \DtN_i \trac_i
   = \bigl(\sum_\mu \lambda_{i,\mu}\respvel_i^{(\mu)}\bigr)^\top
     A_i \bigl(\sum_{\mu'}\lambda_{i,\mu'}\respvel_i^{(\mu')}\bigr)
   = \respvel_i^\top A_i \respvel_i \ge 0$.
\end{proof}

% ----------------------------------------------------------------------

\begin{lemma}[Kernel of the local compliance map]\label{lem:kernel}
  Let $\one_{(0,\normal)} \in \R^{n_i}$ be the vector with entry $1$ at
  every uniform-normal mode $(\gamma,0,\normal)$ (internal interfaces and
  boundary ports alike) and entry $0$ elsewhere.
  Let $L_i \in \R^{n_{i,u} \times n_i}$ be the matrix whose columns are
  all primitive load vectors $\loadvec_i^{(\mu)}$, so that
  $\DtN_i = L_i^\top H_i L_i$.
  If
  \begin{equation}\label{eq:rank-assumption}
    \operatorname{rank}(H_i L_i) = n_i - 1,
  \end{equation}
  then $\ker(\DtN_i) = \operatorname{span}\{\one_{(0,\normal)}\}$ is exactly
  one-dimensional.
\end{lemma}

\begin{proof}
  \textbf{(i)~$\one_{(0,\normal)} \in \ker(\DtN_i)$.}
  Set all uniform-normal coefficients to $p_{\text{const}}$ and all others
  to zero.  Summing the loads:
  \begin{align*}
    \sum_{\gamma} \loadvec_i^{(\gamma,0,\normal)}
    &= \sum_{\gamma} \int_{\iface_{i,\gamma}}
       \bm{\phi}_j^{(i)} \cdot \normal_{i,\gamma} \,\mathrm{d}S
     = \int_{\partial\pore{i}} \bm{\phi}_j^{(i)} \cdot
       \normal \,\mathrm{d}S \\
    &= \int_{\pore{i}} \nabla\cdot\bm{\phi}_j^{(i)} \,\mathrm{d}\bm{x}
     = \bigl[B_i^\top \mathbf{1}\bigr]_j,
  \end{align*}
  where $\mathbf{1}$ is the coefficient vector of $p_h \equiv 1$.
  Hence $(\respvel_i,\resppres_i) = (\bm{0}, p_{\text{const}}\mathbf{1})$
  solves~\eqref{eq:local-sys}, and by uniqueness $\respvel_i = \bm{0}$.
  The energy identity gives $\trac_i^\top \DtN_i \trac_i = 0$, so
  $\DtN_i \trac_i = \bm{0}$ by positive semi-definiteness.

  \textbf{(ii)~$\ker(\DtN_i) \subseteq \operatorname{span}\{\one_{(0,\normal)}\}$.}
  Since $\DtN_i = L_i^\top H_i L_i$ with $H_i$ symmetric,
  $\ker(\DtN_i) = \ker(H_i L_i)$ because $H_i$ is SPD on
  $\operatorname{Range}(L_i)$.  From~\eqref{eq:rank-assumption},
  $\dim\ker(H_i L_i) = n_i - \operatorname{rank}(H_i L_i) = 1$.
  Part~(i) shows $\one_{(0,\normal)} \in \ker(H_i L_i)$, hence
  $\ker(\DtN_i) = \operatorname{span}\{\one_{(0,\normal)}\}$.
  \hfill$\square$
\end{proof}

\begin{remark}\label{rem:rank}
  Assumption~\eqref{eq:rank-assumption} states that the $n_i$ primitive
  responses $H_i \loadvec_i^{(\mu)}$ span a subspace of codimension~$1$ in
  $\operatorname{Range}(H_i)$, the sole dependency being the constant-pressure
  relation $\sum_\gamma \loadvec_i^{(\gamma,0,\normal)} = B_i^\top\mathbf{1}$.
  This is verified numerically for all geometries tested (the local compliance
  matrices have exactly one zero eigenvalue to machine precision) and holds
  analytically provided the nine modes per interface are linearly independent
  on each face and no accidental linear dependencies arise from the load
  assembly.  A rigorous verification for general Lipschitz pore geometries
  requires further work; the numerical evidence is uniform across all tested
  configurations.
\end{remark}

% ----------------------------------------------------------------------

\begin{theorem}[Solvability and mass conservation of Affine-DDPNM]
  \label{thm:main}
  Let $\Schur \in \R^{9m \times 9m}$ be defined by~\eqref{eq:schur}
  and assume Lemma~\ref{lem:kernel} holds for every pore.
  Further assume that every connected component of the pore adjacency
  graph contains at least one pore with a prescribed-traction
  (inlet/outlet) boundary port (the \emph{anchoring condition}).
  Then:
  \begin{enumerate}[label=(\roman*),nosep]
    \item $\Schur$ is symmetric positive-definite (SPD).
    \item The system $\Schur\,\trac = \bm{F}$ has a unique solution.
    \item With $\varepsilon = 0$, the reconstructed velocity field
      satisfies local and global discrete mass conservation
      $\sum_r M_{i,r}^{(0,\normal)} = 0$ (each pore) and
      $\int_{\partial\dom} \vel_h^{\text{DD}}\cdot\normal\,\mathrm{d}S=0$.
  \end{enumerate}
\end{theorem}

\begin{proof}
  \textbf{(i)~SPD.}  Symmetry follows from Lemma~\ref{lem:sym}.
  For any $\trac \in \R^{9m}$, let $\trac_i^u = R_i^u \trac$.
  By~\eqref{eq:energy},
  $\trac^\top \Schur \trac = \sum_i (\trac_i^u)^\top \DtN_i^{uu} \trac_i^u
   = \sum_i \respvel_i^\top A_i \respvel_i \ge 0$.

  If $\trac^\top \Schur \trac = 0$, then $\respvel_i = \bm{0}$ for every
  $i$.  By Lemma~\ref{lem:kernel}, each $\trac_i^u$ equals
  $c_i \one_{(0,\normal)}$ for some $c_i \in \R$, restricted to the
  internal interfaces of $\pore{i}$.  Continuity of $\trac$ across
  interfaces forces $c_i = c$ for all pores in the same connected component
  of the pore graph.  For a component containing a prescribed-traction port,
  the boundary entry of $\one_{(0,\normal)}$ is \emph{not} included in
  $\trac$ (boundary ports belong to $\trac^k$, not $\trac$); hence $c=0$
  on that component.  The anchoring condition ensures every component
  contains a boundary port, so $c=0$ globally.  Hence $\trac = \bm{0}$,
  proving $\Schur$ is SPD.

  \textbf{(ii)} Immediate from (i).

  \textbf{(iii)~Mass conservation.}  With $\varepsilon = 0$, the discrete
  incompressibility $B_i \respvel_i = \bm{0}$ holds.  Choosing the constant
  test function $q_h \equiv 1 \in Q_{i,h}$ and applying the divergence
  theorem yields $\sum_r M_{i,r}^{(0,\normal)} = 0$.  The Schur assembly
  enforces $M_{i,r}^{(0,\normal)} + M_{j,r}^{(0,\normal)} = 0$ on each
  internal interface; summing over all pores, internal fluxes cancel and
  only boundary terms remain, giving the global balance.
\end{proof}

\begin{remark}[Anchoring in practice]
  The anchoring condition is satisfied for all benchmark geometries
  studied: the uniform and random sphere packings have inlet/outlet
  boundaries cutting through the entire domain, and every pore in the
  watershed partition that lacks a solid-wall facet is merged into its
  neighbour (the \emph{floating-basin merge} procedure).  This merge
  step is not an implementation detail---it is the practical enforcement
  of the anchoring condition and is required for the Schur system to be
  non-singular.
\end{remark}

% ======================================================================
\appendix
\section{Error estimate as a Galerkin projection}
\setcounter{lemma}{0}
\setcounter{theorem}{0}
% ======================================================================

We now characterise the Affine-DDPNM as a Galerkin restriction of a
broken-pressure non-overlapping domain-decomposition (C-DD) formulation.
The reference C-DD system uses \emph{discontinuous} pressure across
subdomains ($Q_h^{\text{br}} = \bigoplus_i Q_{i,h}$), not the global
continuous Taylor--Hood space.  Consequently, the error estimate below
quantifies the distance from this broken-pressure C-DD solution, not
from the monolithic Taylor--Hood FEM.  Numerically the two references
differ by $O(10^{-12})$ on the tested meshes (the exact FE-trace Schur
oracle, a primal mixed Schur complement on the global continuous space,
coincides with the global solve to machine precision), but the Galerkin
projection proof itself only covers the broken-pressure C-DD.

% ----------------------------------------------------------------------
\subsection{C-DD formulation (broken pressure)}

On each internal interface $\iface_{i,r}$, the traction is expanded in the
FE face basis $\{\psi_\ell\}$ (the restriction of the velocity trace basis):
\[
  \bm{t}_{i,r}^{\text{FE}}(\bm{x})
  = \sum_{\ell=1}^{m_{i,r}}
    \bigl(\eta_{i,r}^\ell \normal_{i,r}
          + \tau_{i,r}^{\ell,1} \bm{t}_{i,r}^1
          + \tau_{i,r}^{\ell,2} \bm{t}_{i,r}^2\bigr)\,
    \psi_\ell(\bm{x}),
\]
with $m_{i,r}$ velocity DOFs on $\iface_{i,r}$.
Collect the nodal values into
$\bm{\eta}_{i,r},\bm{\tau}_{i,r}^1,\bm{\tau}_{i,r}^2 \in \R^{m_{i,r}}$,
and the full C-DD traction on $\iface_{i,r}$ into
$\bm{\xi}_{i,r} := (\bm{\eta}_{i,r}; \bm{\tau}_{i,r}^1; \bm{\tau}_{i,r}^2)
\in \R^{3m_{i,r}}$.

The Neumann coupling blocks are
\begin{align*}
  N_{i,r}^{\normal} &\in \R^{m_{i,r} \times n_{i,u}},\qquad
  [N_{i,r}^{\normal}]_{\ell,j}
  = \int_{\iface_{i,r}} \psi_\ell\,(\bm{\phi}_j^{(i)}
    \cdot \normal_{i,r})\,\mathrm{d}S,\\
  N_{i,r}^{\bm{t}_\nu} &\in \R^{m_{i,r} \times n_{i,u}},\qquad
  [N_{i,r}^{\bm{t}_\nu}]_{\ell,j}
  = \int_{\iface_{i,r}} \psi_\ell\,(\bm{\phi}_j^{(i)}
    \cdot \bm{t}_{i,r}^\nu)\,\mathrm{d}S \quad (\nu = 1,2).
\end{align*}
Stacking all interfaces vertically gives the local C-DD coupling matrix
\begin{equation}\label{eq:N-stack}
  N_i =
  \begin{bmatrix}
    N_i^{\normal} \\ N_i^{\bm{t}_1} \\ N_i^{\bm{t}_2}
  \end{bmatrix}
  \;\in\; \R^{3\sum_r m_{i,r}\,\times\, n_{i,u}},
\end{equation}
and the local traction $\bm{\xi}_i
= (\bm{\eta}_i; \bm{\tau}_i^1; \bm{\tau}_i^2)$.

The discrete Stokes system on $\pore{i}$ reads
\[
  \begin{bmatrix} A_i & B_i^\top \\ B_i & 0 \end{bmatrix}
  \begin{bmatrix} V_i \\ P_i \end{bmatrix}
  = -\begin{bmatrix} N_i^\top \bm{\xi}_i \\ \bm{0} \end{bmatrix}
    + \text{(boundary contributions)}.
\]
Eliminating interior DOFs yields the local C-DD Schur complement
$\widehat{\Schur}_i = N_i H_i N_i^\top$.
Global Boolean assembly with maps $\widehat{R}_i^u$ gives
\begin{equation}\label{eq:cdd-global}
  \widehat{\Schur}\,\bm{\xi}^{\text{FE}} = \widehat{\bm{F}},\qquad
  \widehat{\Schur} = \sum_{i=1}^N (\widehat{R}_i^u)^\top
    \widehat{\Schur}_i \widehat{R}_i^u,
\end{equation}
with $\bm{\xi}^{\text{FE}} \in \R^{m_\Gamma}$ collecting all global
interface nodal tractions.

% ----------------------------------------------------------------------
\subsection{Affine restriction and equivalence}

For each interface $\iface_k$, let $M_k \in \R^{m_k \times m_k}$ be the
face mass matrix.  For each mode $(\alpha,\beta) \in \modeset$, define
$\bm{b}_k^{\alpha,\beta} \in \R^{m_k}$ by
$(\bm{b}_k^{\alpha,\beta})_\ell
 = \int_{\iface_k} \varphi_\alpha(s,t)\,\psi_\ell\,\mathrm{d}S$.
The per-mode $L^2$-projection vector is
$P_k^{\alpha,\beta} = M_k^{-1} \bm{b}_k^{\alpha,\beta}
 \in \R^{m_k}$.

The local restriction matrix $P_k^{\text{aff}} \in \R^{3m_k \times 9}$
maps the nine affine coefficients to the three-component nodal traction:
\begin{equation}\label{eq:restriction-def}
  P_k^{\text{aff}} =
  \begin{bmatrix}
    P_k^{0,\normal} & P_k^{0,\bm{t}_1} & P_k^{0,\bm{t}_2} &
    P_k^{s,\normal} & P_k^{s,\bm{t}_1} & P_k^{s,\bm{t}_2} &
    P_k^{t,\normal} & P_k^{t,\bm{t}_1} & P_k^{t,\bm{t}_2} \\
    \bm{0} & P_k^{0,\bm{t}_1} & \bm{0} &
    \bm{0} & P_k^{s,\bm{t}_1} & \bm{0} &
    \bm{0} & P_k^{t,\bm{t}_1} & \bm{0} \\[2pt]
    \bm{0} & \bm{0} & P_k^{0,\bm{t}_2} &
    \bm{0} & \bm{0} & P_k^{s,\bm{t}_2} &
    \bm{0} & \bm{0} & P_k^{t,\bm{t}_2}
  \end{bmatrix}.
\end{equation}
Each column occupies only the component block corresponding to its
direction $\beta$:
\begin{itemize}[nosep,leftmargin=*]
  \item $\beta = \normal$: only the first (normal) block-row is non-zero;
  \item $\beta = \bm{t}_1$: only the second block-row is non-zero;
  \item $\beta = \bm{t}_2$: only the third block-row is non-zero.
\end{itemize}
The global restriction operator is the block-diagonal assembly
$\mathcal{P}_{\text{aff}} = \bigoplus_{k=1}^m P_k^{\text{aff}}
: \R^{9m} \to \R^{m_\Gamma}$.

\begin{lemma}[Equivalence]\label{lem:equivalence}
  Let $\Schur$ be the assembled Schur complement~\eqref{eq:schur} and
  $\mathcal{P}_{\text{aff}}$ be defined as above.  Then, for the internal
  degrees of freedom,
  \begin{equation}\label{eq:equivalence}
    \Schur = \mathcal{P}_{\text{aff}}^\top \widehat{\Schur}\,
             \mathcal{P}_{\text{aff}},\qquad
    \bm{F} = \mathcal{P}_{\text{aff}}^\top \widehat{\bm{F}},
  \end{equation}
  where $\widehat{\Schur}$ and $\widehat{\bm{F}}$ are the C-DD system~\eqref{eq:cdd-global}.
\end{lemma}

\begin{proof}
  Fix a pore $\pore{i}$.  For each internal interface $\iface_{i,r}$ and
  each mode $(\alpha,\beta)$, the primitive load vector~\eqref{eq:load}
  equals (up to the sign absorbed in $\bm{d}_{i,\gamma,\beta}$)
  \[
    \loadvec_i^{(\gamma,\alpha,\beta)}
    = -(N_{i,r}^{\beta})^\top P_{i,r}^{\alpha,\beta},
  \]
  where $N_{i,r}^{\beta}$ is the coupling block for direction $\beta$.
  Note that only the $\beta$-component of $P_{i,r}^{\text{aff}}$ is
  non-zero, so the sum over all three directions in the original
  erroneous formula reduces to the single term above.
  The matrix of all $9m_i^u$ primitive internal-interface loads is
  \[
    L_i^u = -(N_i^u)^\top P_i^{\text{aff}},
  \]
  where $N_i^u$ is the C-DD coupling matrix restricted to internal
  interfaces and $P_i^{\text{aff}}$ is the block-diagonal assembly of
  $\{P_{i,r}^{\text{aff}}\}_{r=1}^{m_i^u}$.  The local compliance block is
  $\DtN_i^{uu} = (L_i^u)^\top H_i L_i^u
   = (P_i^{\text{aff}})^\top N_i^u H_i (N_i^u)^\top P_i^{\text{aff}}
   = (P_i^{\text{aff}})^\top \widehat{\Schur}_i^{uu} P_i^{\text{aff}}$.

  Global Boolean assembly with $R_i^u$ (mapping local internal modes to
  $\R^{9m}$) and the C-DD Boolean maps $\widehat{R}_i^u$ (mapping local
  C-DD nodal DOFs to $\R^{m_\Gamma}$) respects the commutative relation
  $\widehat{R}_i^u P_i^{\text{aff}}
   = \mathcal{P}_{\text{aff}}|_{\pore{i}} R_i^u$
  (both sides map the same global affine DOFs to the same global C-DD
  DOFs).  Summing over all pores yields~\eqref{eq:equivalence}.
  The boundary forcing term is treated analogously with prescribed
  $\trac^k$ replacing $\trac$.
\end{proof}

% ----------------------------------------------------------------------
\subsection{Error estimate}

\begin{theorem}[Error estimate for Affine-DDPNM]\label{thm:error}
  Let $\bm{\xi}^{\text{FE}}$ solve the C-DD system~\eqref{eq:cdd-global}.
  Let $\trac^{\text{aff}}$ solve $\Schur\,\trac^{\text{aff}} = \bm{F}$ and
  set $\bm{\xi}^{\text{aff}} = \mathcal{P}_{\text{aff}}\,\trac^{\text{aff}}$.
  Then:
  \begin{enumerate}[label=(\roman*),nosep]
    \item $\|\bm{\xi}^{\text{FE}} - \bm{\xi}^{\text{aff}}\|_{\widehat{\Schur}}
       = \min_{\bm{\mu}\in\R^{9m}}
         \|\bm{\xi}^{\text{FE}}
           - \mathcal{P}_{\text{aff}}\,\bm{\mu}\|_{\widehat{\Schur}}$\label{eq:err-est};
    \item $\sum_i \|\respvel_i^{\text{FE}}
           - \respvel_i^{\text{aff}}\|_{A_i}^2
       = \|\bm{\xi}^{\text{FE}}
         - \bm{\xi}^{\text{aff}}\|_{\widehat{\Schur}}^2$.
  \end{enumerate}
\end{theorem}

\begin{proof}
  By Lemma~\ref{lem:equivalence}, the Affine-DDPNM system is the Galerkin
  restriction of the C-DD system onto $\operatorname{Range}(\mathcal{P}_{\text{aff}})$.
  Galerkin orthogonality gives
  $\mathcal{P}_{\text{aff}}^\top \widehat{\Schur}
   (\bm{\xi}^{\text{aff}} - \bm{\xi}^{\text{FE}}) = \bm{0}$.
  Hence $\bm{e} := \bm{\xi}^{\text{FE}} - \bm{\xi}^{\text{aff}}$ is
  $\widehat{\Schur}$-orthogonal to $\operatorname{Range}(\mathcal{P}_{\text{aff}})$.
  For any $\bm{\mu} \in \R^{9m}$, the Pythagorean identity in the
  $\widehat{\Schur}$-inner product yields
  \[
    \|\bm{\xi}^{\text{FE}} - \mathcal{P}_{\text{aff}}\,\bm{\mu}
      \|_{\widehat{\Schur}}^2
    = \|\bm{e}\|_{\widehat{\Schur}}^2
      + \|\bm{\xi}^{\text{aff}}
        - \mathcal{P}_{\text{aff}}\,\bm{\mu}\|_{\widehat{\Schur}}^2,
  \]
  with equality at $\bm{\mu} = \trac^{\text{aff}}$ (the choice
  $\bm{\mu} = \trac^{\text{aff}}$ makes the last term zero by definition
  of $\bm{\xi}^{\text{aff}}$).
  This proves (i).  Part (ii) follows from the C-DD energy identity
  $\|\bm{\xi} - \tilde{\bm{\xi}}\|_{\widehat{\Schur}}^2
   = \sum_i \|\respvel_i(\bm{\xi}) - \respvel_i(\tilde{\bm{\xi}})\|_{A_i}^2$.
\end{proof}

% ----------------------------------------------------------------------
\subsection{Discussion}

\begin{remark}[Enrichment and energy norm]\label{rem:enrichment}
  The classic DD-PNM restriction $\mathcal{P}_{\text{P0}} : \R^m \to
  \R^{m_\Gamma}$ $L^2$-projects the constant normal function onto the
  FE face basis (normal component only).  Because
  $\operatorname{Range}(\mathcal{P}_{\text{P0}})
   \subset \operatorname{Range}(\mathcal{P}_{\text{aff}})$,
  Theorem~\ref{thm:error} guarantees
  $\|\bm{\xi}^{\text{FE}} - \bm{\xi}^{\text{aff}}\|_{\widehat{\Schur}}
   \le \|\bm{\xi}^{\text{FE}}
     - \bm{\xi}^{\text{P0}}\|_{\widehat{\Schur}}$:
  the Affine error in the C-DD energy norm never exceeds the classic error.
  This monotonicity holds for the $\widehat{\Schur}$-energy norm; it does
  \emph{not} automatically extend to the velocity $L^2$, pressure $L^2$,
  broken $H^1$, or outlet flux errors---those depend additionally on the
  specific geometry and on how the energy norm relates to $L^2$-type norms.
  The 73--90\% $L^2$ reductions observed numerically are \emph{consistent}
  with the energy-norm enrichment guarantee, but are not formally implied
  by it alone.
\end{remark}

\begin{remark}[Mesh dependence]\label{rem:mesh}
  The optimality result~\eqref{eq:err-est} bounds the discrete error
  $\|\bm{\xi}^{\text{FE}} - \bm{\xi}^{\text{aff}}\|_{\widehat{\Schur}_h}$
  in the mesh-dependent energy norm.  A continuum limit argument would
  require: (i)~a fixed interface topology (number of pores and interfaces
  independent of $h$), (ii)~shape-regular mesh refinement with stable
  trace spaces, and (iii)~convergence of the discrete
  Steklov--Poincar\'e operator.  The watershed-based meshes in the
  numerical study do \emph{not} satisfy condition~(i): the basin count
  depends on the mesh size.  The observed $h$-insensitivity of the
  Affine-DDPNM errors within the tested range ($4.3$K--$45$K tetrahedra)
  is therefore a numerical observation, not a consequence of a limiting
  argument.  A rigorous $h$-convergence theory requires a fixed interface
  topology.
\end{remark}

\begin{remark}[Relation to monolithic FEM]
  The broken-pressure C-DD formulation used here differs from the
  monolithic continuous-$P_1$ Taylor--Hood FEM in its pressure space
  ($Q_h^{\text{br}} = \bigoplus_i Q_{i,h}$ vs.\ $Q_h^c \subset H^1$).
  For the geometries tested, the two references coincide numerically to
  $O(10^{-12})$ (verified via the exact FE-trace Schur oracle).  The
  error estimate of Theorem~\ref{thm:error} is therefore a reliable proxy
  for the monolithic FEM error in practice, but a formal proof of
  equivalence between the two reference formulations is beyond the
  present scope.
\end{remark}

% ======================================================================
\begin{thebibliography}{99}

\bibitem{sun2025ddpnm}
S.~Sun, Z.~Wang, L.~Zhang, J.~Zhao.
\newblock A domain decomposition approach to pore-network modeling of porous
  media flow.
\newblock \textit{arXiv:2510.13429v2}, 2026.

\bibitem{horn2012}
R.~A.~Horn and C.~R.~Johnson.
\newblock \textit{Matrix Analysis}, 2nd ed.
\newblock Cambridge University Press, 2012.

\end{thebibliography}

\end{document}
